\documentclass[../main.tex]{subfiles}

\graphicspath{{\subfix{../images/}}}



\begin{document}


\clearpage
\thispagestyle{empty}
\vspace*{7cm}
\section{Introduzione generale}
\vspace*{1.5cm}



\subsection{Introduzione al corso}

Iniziamo ora la trattazione di questo corso andando prima di tutto ad introdurne la sua struttura generale specificando la suddivisione tra lezioni laboratori e completando il tutto con una panoramica sulle modalità d'esame.\\


\subsection{Lezioni}




\subsection{Laboratori}

I laboratori di questo corso, in totale sei, si svolgeranno a settimane alterne a partire dalla terza settimana del corso. I laboratori avranno lo scopo di valutare gli studenti su ciò che é stato spiegato a lezione nelle settimane precedenti.\\
Ad ogni laboratorio si alternerà una una sessione code review all'interno della quale si effettuerà la revisione del codice di un collega al fine di comprendere come gli altri sviluppino e magari trarne dei vantaggi vedendo errori comuni o metodi alternativi di soluzione.\\
Le sessioni di code review possono portare fino ad un massimo di due punti aggiuntivi alla valutazione dell'esame.\\
Si ricorda infine che i laboratori si svolgeranno con appositi strumenti e metodologie introdotti nel corso che saranno fondamentali per affrontare l'esame.\\

\subsection{Esame}

Passiamo infine alla sezione riguardante l'esame ed il suo svolgimento.






	
\end{document}